LLM-based code repair systems iteratively refine generated code using feedback from test execution and static analysis. While prior work has studied which feedback types improve repair quality, the effect of feedback \emph{ordering} under matched content remains unexplored. We hypothesize that presenting static analysis errors before execution failures creates a coarse-to-fine repair trajectory that improves outcomes. Our experiments on HumanEval and MBPP with GPT-4o-mini confirm this: static$\rightarrow$execution ordering achieves 29\% relative improvement over execution$\rightarrow$static ordering, with byte-identical feedback content in both conditions. Investigating the mechanism, we find the effect operates through regression prevention---static-first repairs exhibit 38\% fewer inter-iteration regressions---rather than early-gain acceleration. These findings suggest that feedback structure, not just content, is a first-order design variable in LLM repair systems, offering a free performance gain to practitioners who reorder existing feedback pipelines.
